% Chapitre « Données » (axe 4 et sources communes). Sources : v2/docs/lab/reports/2026-11-14-axe4-*.md,
% 2026-11-14-reconciliation-dentistes.md, 2026-11-14-audit-sources.md (registre C2),
% v2/papers/axe4-donnees/main.tex, v2/data/processed/MANIFEST.md, data/processed/MANIFEST.md,
% v2/docs/lab/decisions.md, fiches v2/docs/lab/fiches/data-*.md. État relu le 2026-11-15.
\providecommand{\donfic}[1]{\nolinkurl{#1}}
\section{Données}\label{don:sec}\label{sec:donnees-ref}
\begingroup\setlength{\emergencystretch}{3em}% local à ce chapitre

Ce chapitre décrit les sources du jumeau, la séparation entre secteur public et secteur privé, le grain
communal, les rapprochements avec des références externes et les limites des données. Abréviations :
\donfic{R/secteur}, \donfic{R/grain}, \donfic{R/projcom}, \donfic{R/flux}, \donfic{R/zonage},
\donfic{R/registre} pour les rapports \donfic{v2/docs/lab/reports/2026-11-14-axe4-*.md} ;
\donfic{R/dentistes} pour \donfic{2026-11-14-reconciliation-dentistes.md} ; \texttt{Audit} pour
\donfic{2026-11-14-audit-sources.md} ; \texttt{A4} pour l'article \donfic{v2/papers/axe4-donnees/main.tex} ;
\texttt{D} pour \donfic{v2/docs/lab/decisions.md}.

\paragraph{Convention.} \mesure{}~: mesuré par le projet sur une donnée ouverte téléchargée, ou lu dans la
source citée à côté du nombre (statut « sourcé » des fiches). \derive{}~: rapport, différence ou arrondi
calculé par nous à partir de nombres publiés ou mesurés. \suppose{}~: valeur par défaut sans mesure, ou seuil
choisi. Les années, les dates, les millésimes et les identifiants (codes FINESS, numéros d'articles) ne sont pas
étiquetés.

\subsection{Périmètre}\label{don:perimetre}

\paragraph{Professions.} Cinq professions sont projetées : sages-femmes (SF), chirurgiens-dentistes (CD),
pharmaciens (PH), masseurs-kinésithérapeutes (MK) et infirmiers (IDE). La médecine générale est hors champ ;
les médecins spécialistes (cardiologie, dermatologie, ophtalmologie) relèvent d'une phase ultérieure
(\donfic{CLAUDE.md}, recadrage du 2026-11-14 ; \donfic{v2/ROADMAP.md}).

\paragraph{« Libéral » désigne le secteur privé.} La sortie du jumeau s'appelle « libéral » dans l'interface et
\texttt{prive} dans le code. Elle compte tout professionnel dont l'employeur n'est pas public (\texttt{D},
entrée du 2026-11-14) :
\begin{itemize}
  \item \textbf{inclus} : libéraux et mixtes ; pharmaciens adjoints salariés d'officine ; salariés des cliniques
  privées et des établissements privés non lucratifs ; industrie ; laboratoires privés ;
  \item \textbf{exclus} : hôpital public, protection maternelle et infantile (PMI), Éducation nationale,
  EHPAD publics, centres de santé municipaux et tout autre employeur public.
\end{itemize}
Les sources emploient « libéral » au sens d'exercice indépendant (RPPS, data.ameli.fr). Le mot n'a donc pas le
même sens dans le jumeau et dans les fichiers : la sortie du jumeau est plus large que les « libéraux » de
l'Assurance maladie, puisqu'elle inclut les salariés du privé.

\subsection{Catalogue des sources}\label{don:catalogue}

Le tableau~\ref{don:tab-catalogue} liste les jeux utilisés. Toutes les sources sont ouvertes. Les données
brutes ne sont pas versionnées (\donfic{data/raw/}, \donfic{v2/data/raw/}) ; les sorties traitées le sont
(section~\ref{don:manifeste}). Les scripts sont dans \donfic{v2/data/} sauf mention « v1 » (\donfic{v2/sante/}).
Les fichiers produits sont dans \donfic{v2/data/processed/} sauf mention \donfic{data/processed/}.

\paragraph{Licences.} Trois licences seulement sont écrites dans le dépôt :
\begin{itemize}
  \item fiche Atlasanté c987d0e2 : « utilisation libre sous réserve de mentionner la source et la date de
  dernière mise à jour » (\donfic{R/zonage}) ;
  \item jeu DREES n\textsuperscript{o}~491 : Licence Ouverte (bibliographie de l'article de l'axe~1)
  \citep{d-drees491} ;
  \item zonages data.gouv \donfic{t-zonages-com.csv} : Licence Ouverte v2 (même bibliographie)
  \citep{d-zonagesdatagouv}.
\end{itemize}
Pour les autres jeux, aucun texte lu ne nomme la licence : la colonne porte « non vérifiée ». Que ces jeux
publics soient sous Licence Ouverte Etalab est plausible, mais c'est une supposition, à vérifier avant toute
diffusion.

{\scriptsize
\setlength{\tabcolsep}{2pt}
\begin{longtable}{@{}p{1.45cm}p{2.45cm}p{1.6cm}p{2.15cm}p{1.35cm}p{1.3cm}p{2.0cm}p{2.5cm}@{}}
\caption{Catalogue des jeux de données. Producteur, millésime et grain d'après les rapports de l'axe~4, le
registre C2 de l'audit et les manifestes. « Fiches » : \donfic{v2/docs/lab/fiches/}.}\label{don:tab-catalogue}\\
\toprule
Producteur & Jeu & Millésime & Grain & Période & Licence & Script & Fichier produit \\
\midrule\endfirsthead
\toprule
Producteur & Jeu & Millésime & Grain & Période & Licence & Script & Fichier produit \\
\midrule\endhead
\bottomrule\endfoot
\multicolumn{8}{@{}l}{\emph{Stocks nationaux et départementaux}}\\
DREES & Démographie des professionnels de santé depuis 2012 (RPPS ; ADELI pour les MK avant 2017)
\citep{d-drees2026demo} & publication du 2026-07-02 (feuilles Lisez-moi) & dép.\ $\times$ secteur
(8\mesure{} à 12\mesure{} modalités) $\times$ mode (4\mesure) $\times$ âge (10\mesure{} tranches)
$\times$ sexe, au 1\textsuperscript{er} janvier & 2012--2026 ; MK : ADELI 2012--2016, RPPS 2017--2026 &
non vérifiée & \donfic{build_secteur.py} ; v1 : \donfic{load.py} & \donfic{secteur_stock.parquet}
(\num{899250}\mesure{} lignes), \donfic{secteur_split.parquet}, \donfic{secteur_checks.csv},
\donfic{secteur_finess_shares.csv} ; v1 : \donfic{data/processed/stock.parquet}, \donfic{checks.csv} \\
DREES, Insee & Infirmières salariées (Base Tous Salariés, BTS) et libérales (SNDS)
\citep{d-dreesbts2026} & BTS : 2025-03-25, mise à jour 2026-07-02 ; SNDS : 2025-03-25 & dép.\ $\times$
secteur (6\mesure) ; aucun croisement secteur $\times$ âge ou sexe & BTS 2013--2023 ; SNDS 2012--2023 &
non vérifiée & \donfic{build_secteur.py}, \donfic{build_infirmiers.py} ; v1 : \donfic{nurses.py} &
\donfic{infirmiers_secteur_age_sexe.parquet}, \donfic{infirmiers_test_hind.csv} ; v1 :
\donfic{stock_nurses_salaried.parquet} \\
CNAM & data.ameli.fr, démographie des professionnels libéraux & modifié le 2026-10-28 & dép., âge par
5~ans, sexe ; sans les PH & 2010--2025 & non vérifiée & v1 : \donfic{ameli.py} &
\donfic{data/processed/stock_ameli.parquet} \\
Insee & Omphale 2022, 3\mesure{} scénarios (central, population basse, population haute) & 2022 &
dép.\ $\times$ âge $\times$ sexe & projection & non vérifiée & v1 : \donfic{load.py} &
\donfic{data/processed/population.parquet} \\
\midrule
\multicolumn{8}{@{}l}{\emph{Établissements et statut juridique}}\\
Ministère chargé de la santé & FINESS établissements (\texttt{etalab-cs1100502}, \SI{35.8}{\mega\byte}\mesure)
et entités juridiques (\texttt{cs1100501}, \SI{11.9}{\mega\byte}\mesure) \citep{d-finessEt2026,d-finessEj2026}
& arrêté au 2026-05-04, en ligne le 2026-05-12 & établissement, catégorie, statut juridique & instantané 2026
& non vérifiée & \donfic{build_secteur.py} ; \donfic{build_annuaire.py} & \donfic{secteur_finess_shares.csv} \\
idem & FINESS historique (\SI{196}{\mega\byte}\mesure) & état au 31/12 de chaque année & établissement ;
catégories 620, 603 (depuis 2012), 124 (depuis 2014) & 2004--2025 & non vérifiée &
\donfic{build_finess_histo.py} & \donfic{finess_commune_annee.parquet}, \donfic{officines_flux_commune.parquet},
\donfic{finess_histo_checks.csv} \\
\midrule
\multicolumn{8}{@{}l}{\emph{Grain communal}}\\
ANS & Annuaire Santé, extraction libre accès du RPPS (\donfic{PS_LibreAcces.zip},
\SI{228}{\mega\byte}\mesure) \citep{d-ans2026annuaire} & 2026-11-13 & commune de chaque activité
(arrondissements PLM) ; civilité, mode, secteur, FINESS, origine du diplôme ; sans âge & instantané unique &
non vérifiée & \donfic{build_annuaire.py} ; \donfic{build_commune.py} ; \donfic{archive_annuaire.py} &
\donfic{annuaire_commune.parquet}, \donfic{annuaire_checks.csv}, \donfic{annuaire_checks_dep.csv} ; archive
brute non versionnée (\num{813557}\mesure{} activités, \num{752993}\mesure{} personnes) \\
Atlasanté (FNPS, SNDS) & CartoSanté, séries \donfic{*_popage} (service TJS) \citep{d-cartosante2026} &
libéraux actifs au 31/12 & commune COG 2026 et \texttt{arm2026} ; 5\mesure{} tranches d'âge ; sans sexe &
IDE 2004--2025, MK 2005--2025, CD 2006--2025 ; SF rejetée & non vérifiée & \donfic{build_commune.py},
\donfic{build_flux.py} ; v1 : \donfic{cartosante.py} & \donfic{commune_stock.parquet},
\donfic{commune_geo.parquet}, \donfic{commune_checks.csv}, \donfic{grain_petits_effectifs.csv},
\donfic{grain_bruit.csv}, \donfic{flux_installations_commune.parquet}, \donfic{flux_checks.csv} ; v1 :
\donfic{stock_commune.parquet} \\
CNAMTS, ARS & CartoSanté \donfic{zonage_conv} (4\mesure{} fichiers XML d'environ
\SI{1.6}{\mega\byte}\mesure) \citep{d-cartosante2026} & extrait le 2026-11-14 ; date des arrêtés non servie &
\num{34919}\mesure{} géocodes & actuel, non daté & non vérifiée & \donfic{build_zonage.py} &
\donfic{zonage_commune.parquet} \\
ARS Aquitaine (DOSA) & Atlasanté, fiche c987d0e2 (GeoJSON, \SI{6.4}{\mega\byte}\mesure)
\citep{d-atlasante2014} & publiée le 2014-01-22, révisée le 2014-10-22 & \num{36700}\mesure{} communes,
métropole ; IDE, MK, CD & génération « 2012 », déduite & utilisation libre, source et date à mentionner &
\donfic{build_zonage.py} & \donfic{zonage_commune.parquet} (colonne \texttt{millesime}) \\
ARS, CNAM (Atlasanté, DataSanté) & Zonages des professionnels de santé libéraux, \donfic{t-zonages-com.csv}
\citep{d-zonagesdatagouv} & ressource du 2026-10-30, zonage au 2026-10-20 & commune, niveaux 1 à 5 & actuel &
Licence Ouverte v2 & aucun (fiches \emph{$\ast$-regul-regle}) & aucun \\
Insee & Bassins de vie 2022 au 2026-01-01 \citep{d-insee2026bv} ; COG 2026 et événements depuis 1943
\citep{d-insee2026cog} ; séries historiques de population 1876--2023 ; recensements 2012, 2017, 2023 ; grille
de densité ; aires d'attraction des villes 2020 & selon le jeu & commune & selon le jeu & non vérifiée &
\donfic{build_commune.py}, \donfic{build_passage.py}, \donfic{build_attributs.py} &
\donfic{passage_commune.parquet} (\num{39227}\mesure{} codes), \donfic{petits_effectifs.parquet},
\donfic{petits_effectifs_classes.csv}, \donfic{commune_attributs.parquet} (\num{490154}\mesure{} lignes),
\donfic{commune_attributs_checks.csv} \\
Etalab & Découpage administratif, geo.api.gouv.fr \citep{d-geoapi2026} & lu le 2026-11-14 & commune
(\num{34969}\mesure), arrondissement (45\mesure), EPCI (\num{1255}\mesure) & 2026 & non vérifiée & contrôles
& aucun \\
Projet (dérivé) & Pharmaciens par commune reconstitués (DREES, FINESS historique, Annuaire) & --- & commune
$\times$ an & 2012--2026 & --- & \donfic{build_pharmaciens.py} & \donfic{pharmaciens_commune_annee.parquet},
\donfic{pharmaciens_checks.csv}, \donfic{pharmaciens_checks_dep.csv} \\
\midrule
\multicolumn{8}{@{}l}{\emph{Formation, établissements, événements (sources des paramètres)}}\\
DREES & Enquête Écoles, jeu 491, séries longues \citep{d-drees491} & mise à jour du 2026-10-29 & national &
1966--2025 & Licence Ouverte & script en annexe de l'article de l'axe~1 & fiches \emph{ide-$\ast$} \\
DREES & SAE, bases statistiques, jeu 708 \citep{d-sae708} & mise à jour du 2026-11-04 & établissement &
2023--2025 & non vérifiée & script de la fiche \emph{part-privee-salaries-sae} & fiche \\
DREES & Personnel des services de PMI \citep{d-dreespmi} & au 31/12/2024 & national & 2010--2024 & non
vérifiée & aucun & override SF de \donfic{build_secteur.py} \\
Projet & Registre des réformes : Légifrance (via reproductions JO), ameli, Conseil d'État, DREES, ARS, Ordres &
--- & événement daté & 2009--2028 & --- & versionné tel quel & \donfic{v2/data/events.csv}
(30\mesure{} lignes) \\
\end{longtable}
}

\subsection{Séparation public / privé}\label{don:secteur}

\paragraph{Pourquoi.} L'ancienne sortie « libéral + mixte » ignorait les salariés du privé. L'ancienne
fourchette des SF 2026 allait de \num{12475}\mesure{} à \num{26428}\mesure{} (\donfic{R/secteur}). Les fichiers
DREES croisent pourtant le secteur d'activité avec le département, le mode, l'âge et le sexe
\citep{d-drees2026demo} : la séparation était déjà dans les données (\texttt{A4}, \S~Méthode).

\paragraph{Méthode (\donfic{build_secteur.py}).} Chaque cellule (profession, secteur, mode, département, âge,
sexe, année) reçoit une part privée $\pi$ :
\begin{itemize}
  \item libéral ou mixte : $\pi = 1$\derive{}, quel que soit le secteur ;
  \item salarié d'un secteur sans ambiguïté : $\pi = 0$\derive{} à l'hôpital public ; $\pi = 1$\derive{} en
  ESPIC, privé lucratif, officine, laboratoire, industrie et cabinet ;
  \item salarié d'un secteur mixte (centre de santé, prévention, médico-social, « autre ») :
  $\pi \in [0 ; 1]$, avec une valeur centrale $\hat\pi$.
\end{itemize}
Trois sorties en découlent : $N_{\min}$, somme des cellules sûrement privées ; $N_{\max} = N_{\min}$ + somme des
cellules ambiguës ; $\hat N = \sum \hat\pi\, N$.

\paragraph{Parts centrales par FINESS.} $\hat\pi$ est la part des établissements d'une catégorie FINESS dont
l'entité juridique est privée. Les statuts juridiques 1 à 46 (sécurité sociale comprise) sont comptés publics,
les statuts 47 et plus privés (\donfic{R/secteur}). Parts mesurées (tableau~\ref{don:tab-finess}).

\begin{table}[htbp]
\centering\small
\caption{Parts privées par catégorie FINESS, 2026\mesure{} (\donfic{R/secteur}, \donfic{secteur_finess_shares.csv}).}
\label{don:tab-finess}
\begin{tabular}{@{}lrr@{}}
\toprule
Catégorie FINESS & Part privée & Établissements \\
\midrule
Centres de santé (124) & \num{0.798}\mesure & \num{3395}\mesure \\
\quad dont centres « dentaires » & \num{0.963}\mesure & \num{1299}\mesure \\
PMI et planification (223, 228) & \num{0.070}\mesure & \num{1435}\mesure \\
EHPAD (500) & \num{0.558}\mesure & \num{7399}\mesure \\
Handicap médicalisé & \num{0.866}\mesure & \num{6008}\mesure \\
Médico-social, agrégat appliqué au secteur « Établissement sanitaire et social » des MK, des SF et des CD & \num{0.696}\mesure & \num{13407}\mesure \\
\bottomrule
\end{tabular}
\end{table}

Deux secteurs n'ont pas de proxy FINESS : « Autre secteur », et « Prévention et établissement de santé » chez
les PH. Leur part vaut $\hat\pi = 0{,}5$\suppose{}. Cela fait 5\mesure{} cellules profession $\times$ secteur
(PH autre, PH prévention, MK autre, SF autre, CD autre ; \texttt{Audit}, registre C2, ligne « trois parts »
corrigée).

\paragraph{Exception des SF en « Prévention ».} Le 2026-11-14 au soir, les SF salariées du secteur
« Prévention » ont été classées publiques ($\pi = 0$\derive{}, constante \texttt{OVERRIDE} de
\donfic{build_secteur.py} ; \texttt{D}, entrée « Axe 4 : SF \enquote{Prévention} = public »). Motif : la DREES
compte \num{1423}\mesure{} SF en PMI au 31/12/2024 \citep{d-dreespmi}, et la PMI est un service
départemental, donc un employeur public.

\paragraph{Infirmiers.} Le BTS sépare déjà public et privé en trois paires (hôpital, médico-social, autre), et
les libéraux du SNDS sont privés \citep{d-dreesbts2026} : il n'y a pas d'ambiguïté. Mais le BTS ne donne ni
âge ni sexe par secteur. La ventilation suppose $H_{\mathrm{ind}}$ : dans un département, le secteur est
indépendant de l'âge et du sexe. Le test sur le sexe, au seuil de 3\suppose{} points fixé avant calcul, donne
un écart signé médian de \num{2.1}\mesure{} points et un écart absolu médian de \num{2.45}\mesure{} points, sous le seuil de 3 : l'hypothèse n'est pas rejetée (\donfic{infirmiers_test_hind.csv}, 104\mesure{} départements).
Elle n'est pas testée sur l'âge.

\paragraph{Fourchettes.} Le tableau~\ref{don:tab-fourchettes} donne la version en vigueur.

\begin{table}[htbp]
\centering\footnotesize
\setlength{\tabcolsep}{2.5pt}
\caption{Effectif « libéral » (= privé), France entière, dernière année disponible (\texttt{A4},
tab.~2 ; \texttt{D}, entrée du 2026-11-14 au soir). Largeur : $N_{\max}-N_{\min}$, avant et après la séparation.}
\label{don:tab-fourchettes}
\begin{tabular}{@{}lrrrrrrr@{}}
\toprule
Prof. & Année & Tous secteurs & Ancienne fourchette & Min & Centrale & Max & Largeur \\
\midrule
SF & 2026 & \num{26428}\mesure & \num{12475}--\num{26428}\mesure & \num{12059}\mesure & \num{12819}\derive & \num{13331}\mesure & \num{13953} $\to$ \num{1272}\derive \\
PH & 2026 & \num{75335}\mesure & \num{65344}--\num{75335}\mesure & \num{60005}\mesure & \num{63612}\derive & \num{67219}\mesure & \num{9991} $\to$ \num{7214}\derive \\
CD & 2026 & \num{48690}\mesure & \num{48143}--\num{48690}\mesure & \num{40784}\mesure & \num{47349}\derive & \num{47894}\mesure & \num{547} $\to$ \num{7110}\derive \\
MK & 2026 & \num{111141}\mesure & \num{98348}--\num{111237}\mesure & \num{98828}\mesure & \num{101951}\derive & \num{103986}\mesure & \num{12889} $\to$ \num{5158}\derive \\
IDE & 2023 & \num{607644}\mesure & --- & \num{277788}\mesure & \num{277788}\mesure & \num{277788}\mesure & 0\mesure \\
\bottomrule
\end{tabular}
\end{table}

La largeur baisse de 91~\%\derive{} pour les SF, 60~\%\derive{} pour les MK et 28~\%\derive{} pour les PH.
Elle augmente pour les CD : les anciennes bornes basses n'étaient pas des minima, car elles comptaient tout
« autre salarié » comme privé, PMI et centres municipaux compris ; les \num{6694}\mesure{} CD salariés des
centres de santé sont désormais ambigus, avec $\hat\pi = \num{0.963}$\mesure. Infirmiers 2023 : \num{98854}\mesure{}
libéraux (SNDS) et \num{178934}\mesure{} salariés privés (BTS : hôpital \num{94347}, médico-social \num{40050},
autre \num{44537}\mesure) ; public \num{329856}\mesure{} (\texttt{Audit}, registre C2).

\paragraph{L'incohérence sur les SF, tranchée.} Deux fourchettes SF 2026 circulent :
\begin{itemize}
  \item \num{12059}--\num{14660}\mesure{} (centrale \num{12912}, largeur $-81$~\%\derive) dans le
  \texttt{headline} et les tableaux de \donfic{R/secteur} ;
  \item \num{12059}--\num{13331}\mesure{} (centrale \num{12819}\derive, largeur $-91$~\%\derive) dans la mise à
  jour du soir de ce même rapport, dans \texttt{A4}, la fiche \emph{data-secteur-sf} et \texttt{D}.
\end{itemize}
\textbf{La valeur en vigueur est \num{12059}--\num{13331}.} La première est le calcul d'avant l'exception
« Prévention ». Dans ce premier calcul, les \num{1329}\mesure{} SF salariées en « Prévention » étaient ambiguës,
avec $\hat\pi = \num{0.070}$\mesure{} (part privée FINESS des PMI et centres de planification). Les classer
publiques les retire du maximum et presque rien de la centrale :
\[
\num{14660} - \num{1329} = \num{13331}\derive, \qquad
\num{12912} - \num{0.070}\times\num{1329} \approx \num{12819}\derive .
\]
La largeur passe de \num{2601}\derive{} à \num{1272}\derive{}, soit de $-81$~\% à $-91$~\%\derive{} par
rapport aux \num{13953}\derive{} d'origine. Le rapport signale la mise à jour en tête, mais garde ses anciens
tableaux et son ancien \texttt{headline} : c'est la source de l'écart.

Un second écart en découle : \num{1329}\mesure{} SF « Prévention » dans \donfic{R/secteur}, \num{1383} dans \texttt{A4},
l'article de l'axe~1 et \texttt{D}. Nous l'avons mesuré sur \donfic{secteur_stock.parquet} (2026) : le
secteur « Prévention » compte \num{1329}\mesure{} salariées, classées publiques, et \num{54}\mesure{} mixtes,
privées par la règle libéral ou mixte, soit \num{1383}\mesure{} au total. Les deux nombres sont justes ; ils
portent sur deux champs (salariées seules, tous modes).

\paragraph{Ambiguïté restante.} Les secteurs mixtes de plus de 500\mesure{} personnes en 2026 sont :
PH « Autre secteur » \num{5805}\mesure{} et « Prévention et établissement de santé » \num{1409}\mesure{}
($\hat\pi = 0{,}5$\suppose) ; CD centres de santé \num{6694}\mesure{} ($\hat\pi = \num{0.963}$\mesure) ; MK
établissements sanitaires et sociaux \num{2200}\mesure{} ($\hat\pi = \num{0.696}$\mesure) et « Autre secteur »
\num{2524}\mesure{} ($0{,}5$\suppose) ; SF « Autre secteur » \num{796}\mesure{} ($0{,}5$\suppose), soit
63~\%\derive{} du résidu SF. La largeur médiane de la fourchette par département vaut 3,6~\%\mesure{} de
l'effectif pour les SF, 4~\%\mesure{} pour les MK et 9~\%\mesure{} pour les CD et les PH (\donfic{R/secteur}).

\paragraph{Pharmaciens hospitaliers : trois sources, trois valeurs.} Le RPPS place 16,6~\%\mesure{} (min),
18,2~\%\mesure{} (centrale) et 19,8~\%\mesure{} (max) des PH salariés hospitaliers dans le privé en 2026. La SAE
en donne 29,4~\%\mesure{} \citep{d-sae708}, l'Ordre 28,6~\%\derive{} des inscriptions des sections H et E
\citep{d-cnop2025}. La variante SAE ajouterait \num{1116}\derive{} privés (centrale \num{63612} $\to$
\num{64728}\derive), dans la fourchette. Le jumeau garde la centrale RPPS, seule croisée par département, âge et
sexe (\texttt{D}). L'écart n'est pas réconcilié.

\paragraph{Réserves de méthode.}
\begin{itemize}
  \item Le secteur DREES est celui de l'activité la plus récente, pas de l'activité principale
  \citep{d-drees2026demo} : un salarié public qui prend un second poste privé est compté privé.
  \item Les parts FINESS comptent des établissements, pas des salariés. La pondération par l'effectif n'est
  pas mesurée.
  \item Les parts FINESS 2026 sont appliquées à toutes les années 2012--2026\suppose.
  \item Le SNDS et le BTS peuvent compter deux fois une même infirmière ; ce recouvrement n'est pas mesuré.
\end{itemize}

\subsection{Grain communal et petits effectifs}\label{don:grain}

\paragraph{Trois sources descendent à la commune.} L'Annuaire Santé couvre les cinq professions, mais c'est
un instantané (2026-11-13), sans âge ni date de naissance \citep{d-ans2026annuaire}. CartoSanté donne
l'histoire des libéraux par tranche d'âge pour les IDE, les MK et les CD \citep{d-cartosante2026}. FINESS décrit
les structures. La série SF de CartoSanté (\donfic{sf_popage}) est rejetée : ses modalités d'âge sont libellées
en années (\donfic{R/grain}).

\paragraph{Les effectifs communaux sont petits.} La médiane des effectifs privés par commune est nulle pour
les cinq professions (tableau~\ref{don:tab-grain}).

\begin{table}[htbp]
\centering\small
\caption{Effectifs privés par unité, Annuaire Santé du 2026-11-13\mesure{} (\donfic{R/grain} ;
\donfic{grain_petits_effectifs.csv}). Unités : \num{34917}\mesure{} communes (arrondissements PLM),
\num{1707}\mesure{} bassins de vie, \num{1255}\mesure{} EPCI, 101\mesure{} départements.}
\label{don:tab-grain}
\begin{tabular}{@{}lrrrrrr@{}}
\toprule
 & \multicolumn{4}{c}{Médiane par unité} & \multicolumn{2}{c}{Part des pros dans une unité $<30$} \\
\cmidrule(lr){2-5}\cmidrule(l){6-7}
Profession & Commune & Bassin de vie & EPCI & Dép. & Commune & EPCI \\
\midrule
IDE & 0\mesure & \num{35.5}\mesure & 53\mesure & \num{1644}\mesure & 30,6~\%\mesure & 2,6~\%\mesure \\
MK  & 0\mesure & 17\mesure & 23\mesure & 615\mesure & 49,8~\%\mesure & 10,5~\%\mesure \\
PH  & 0\mesure & 11\mesure & \num{16.5}\mesure & 444\mesure & 57,5~\%\mesure & 18,4~\%\mesure \\
CD  & 0\mesure & 6\mesure & 8\mesure & 261\mesure & 53,4~\%\mesure & 20,0~\%\mesure \\
SF  & 0\mesure & \num{1.5}\mesure & \num{2.5}\mesure & 78\mesure & 69,8~\%\mesure & 39,9~\%\mesure \\
\bottomrule
\end{tabular}
\end{table}

\begin{itemize}
  \item \textbf{Communes sans aucun privé}, sur \num{34917}\mesure{} : IDE \num{22657}\mesure, MK
  \num{25296}\mesure, PH \num{26704}\mesure, CD \num{29538}\mesure, SF \num{31951}\mesure.
  \item \textbf{EPCI de moins de 5 privés} : SF 881\mesure, CD 429\mesure, PH 106\mesure, MK 73\mesure, IDE
  6\mesure{} (\donfic{R/grain} ; \donfic{R/projcom} donne 878, 427, 105, 71 et 5\mesure, calcul de
  \donfic{build_passage.py} ; annexe~\ref{don:ann}).
  \item \textbf{Bruit d'une année sur l'autre} (CartoSanté, médiane de $|\Delta \log N|$ entre 2024 et 2025) :
  5 à 7~\%\mesure{} par an dans la classe 5--29, 1,5 à 3~\%\mesure{} dans la classe 100 et plus. Les flux pèsent
  5 à 10~\%\derive{} du stock par an, soit moins d'une installation par an dans la plupart des communes de 5 à
  29 professionnels\derive{} (\donfic{grain_bruit.csv}).
  \item \textbf{Diffusion} : 52,5~\%\mesure{} des cellules non nulles IDE, MK et CD $\times$ commune $\times$ âge
  $\times$ année (2004--2025) de \donfic{commune_stock.parquet} valent une personne
  (\donfic{2026-11-15-audit-rgpd-aipd-dm.md}) ; en 2025 seule, 49,2~\%\mesure{} des \num{67914}\mesure{}
  cellules non nulles. Le seuil de
  diffusion des tables communales publiées est une question ouverte.
\end{itemize}

\paragraph{Grain projetable.} La règle retenue est une convention : médiane d'au moins 30\suppose{} privés par
unité, soit un bruit de Poisson d'environ $1/\sqrt{30} \approx 18$~\%\derive. Elle donne le bassin de vie ou
l'EPCI pour les IDE ; l'EPCI, avec un modèle hiérarchique, pour les MK et les PH ; le département pour les CD et
les SF (\texttt{A4}, \S~Grain). Seul le banc de validation (protocole N10, axe~3) peut remplacer ce seuil. La
commune reste un grain d'observation et de répartition.

\paragraph{Ce qui a été reconstitué.}
\begin{itemize}
  \item \textbf{Pharmaciens par commune, 2012--2026} : reconstitués à partir de la DREES, de FINESS historique
  et de l'Annuaire ; les parts hors officine sont tenues à leur niveau de 2026\suppose{} (\donfic{R/projcom}).
  \item \textbf{SF} : aucun historique communal ouvert. Cinq routes ont été testées (\donfic{sf_popage},
  \donfic{sf_datcom.actifs}, \donfic{offre_ps.dens_sf}, ameli, Annuaire). Le repli (parts communales 2026
  $\times$ série départementale) est modélisé, non observé.
  \item \textbf{Flux d'installation} : aucune source ouverte par commune. Le substitut $I_{\mathrm{net}}$ est un
  solde (entrées moins sorties des moins de 40~ans) qui compte les déménagements et ne voit pas le salariat
  (\donfic{R/flux}).
\end{itemize}

\subsection{Rapprochements avec des références externes}\label{don:rapprochements}

\begin{itemize}
  \item \textbf{Cellules DREES contre total France du fichier} : 0~\%\mesure{} d'écart pour SF, CD, PH et MK
  ADELI ; MK RPPS 0,22~\%\mesure{} au plus (2026) ; BTS 0,009~\%\mesure{} au plus. L'écart MK est interne au
  fichier : en 2017, \num{84599}\mesure{} (cellules) contre \num{84449}\mesure{} (total France) ; la feuille
  Lisez-moi prévient que certains totaux peuvent différer de la somme \citep{d-drees2026demo}.
  \item \textbf{CartoSanté (au 31/12/$Y$) contre RPPS libéral + mixte et SNDS (même année $Y$)} : MK 2022--2025
  $-3{,}7$, $-3{,}1$, $-2{,}5$, $-3{,}2$~\%\mesure{} ; CD $-3{,}3$, $-2{,}7$, $+0{,}3$, $-0{,}8$~\%\mesure{} ;
  IDE 2020--2023 $+3{,}9$, $+3{,}5$, $+4{,}0$, $+4{,}7$~\%\mesure{} (\num{103499} contre \num{98854}\mesure{} en
  2023). En alignant le 31/12/$Y$ sur le 1/1/$Y+1$ : MK de $-7{,}2$ à $-6{,}0$~\%\mesure{}, CD de $-4{,}4$ à
  $-2{,}8$~\%\mesure{}, IDE de $+2{,}6$ à $+4{,}4$~\%\mesure. L'écart MK connu de $-5{,}4$~\%\mesure{} en 2025
  tombe entre ces deux alignements : il dépend de la date de référence\derive{} (\donfic{R/grain}).
  \item \textbf{Annuaire (privés, novembre 2026) contre centrale DREES 2026, médiane par département} (déciles
  1 et 9) : PH $-3{,}2$~\% ($-9{,}2$ ; $-0{,}3$), MK $-5{,}5$~\% ($-9{,}3$ ; $-0{,}7$), CD $-10{,}8$~\%
  ($-16{,}5$ ; $-5{,}5$), SF $-10{,}8$~\% ($-22{,}3$ ; $-1{,}8$), IDE $-14{,}6$~\% ($-22{,}8$ ;
  $-7{,}8$)\mesure. La référence IDE est 2023, avec trois ans d'écart et un double compte possible. Au national,
  l'Annuaire compte \num{61882}\mesure{} PH privés, dans la fourchette \num{60005}--\num{67219}.
  \item \textbf{Pharmaciens d'officine reconstitués contre Annuaire, 2026} : \num{50196} contre
  \num{50664}\mesure{} ($-0{,}9$~\%\derive) ; corrélation par commune \num{0.98}\mesure{} ; erreur médiane
  16,1~\%\mesure{} à la commune (seuil 25~\%\suppose, passé) et 10,7~\%\mesure{} à l'EPCI (seuil
  10~\%\suppose, échoué de peu) (\donfic{R/projcom}).
  \item \textbf{Officines FINESS contre Ordre} : \num{20324}\mesure{} fin 2024 contre \num{20242}\mesure{} au
  1/1/2025 ($+0{,}4$~\%\derive) \citep{d-ordre2025demo}.
  \item \textbf{Hospitaliers, SAE contre RPPS} : SF 19,8 contre 17,4~\%\mesure{} ; MK 44,3 contre
  44,4~\%\mesure{} ; CD 2,2 contre 4,3~\%\mesure{} (24 privés sur 557 salariés hospitaliers au RPPS 2025, recompté le 15/11 ; le 3,4~\% de la fiche \texttt{param-part-privee-salaries-sae} ne se reproduit pas) \citep{d-sae708}. Concordance à 2\derive{} points pour SF et
  MK ; les PH divergent (section~\ref{don:secteur}).
  \item \textbf{Zonage dentaire} : 97\mesure{} géocodes « non prioritaires » et 4,8~\%\mesure{} de la population
  (\num{3309773} sur \num{68392618}\mesure) ; l'ARS PACA cite 5~\%\mesure{} \citep{d-arspaca2025cd}, le mémoire
  EHESP 97\mesure{} communes (p.~53, confirmé par la relecture du 15/11) \citep{d-bahroun2025}.
  \item \textbf{Zonage CartoSanté contre data.gouv} : parts de population IDE 27,6 contre 27,4~\%, MK 27,4
  contre 26,5~\%, SF 12,3 contre 12,2~\%, CD 4,8~\% dans les deux\mesure{} (\donfic{R/zonage} ;
  \citep{d-zonagesdatagouv}).
  \item \textbf{Année blanche MK 2018} \citep{d-er1159} : les moins de 25~ans passent de \num{5634} à
  \num{3687}\mesure{} au RPPS ($-34{,}6$~\%\derive), puis remontent à \num{5137}\mesure{} en 2026. Le substitut
  $I_{\mathrm{net}}$ ne la voit qu'en partie : \num{2604} contre \num{2877}\mesure{}, rapport
  \num{0.90}\derive{} au-dessus du seuil de \num{0.8}\suppose{} (test T2 échoué ; \donfic{R/flux}).
\end{itemize}

\paragraph{Dentistes : « \num{76000} contre \num{62000} ».} Deux chiffres circulaient dans l'équipe
(\donfic{R/dentistes}). \num{76397}\mesure{} est la valeur 2050 de la projection de référence (API
\donfic{/api/scenario}, \texttt{base[2050]}), pas le socle. Le socle privé 2026 vaut \num{47347.5}\mesure{} hors code département 999 (aucune ligne en code 20) ; la centrale de \donfic{secteur_stock.parquet}, tous codes, vaut \num{47349.5}\mesure{},
sous les \num{48690}\mesure{} tous modes (libéral \num{36259}, mixte \num{3824}, autre salarié \num{8060},
salarié hospitalier \num{547}\mesure). Aucun double compte : 0\mesure{} clé dupliquée sur \num{13909}\mesure{}
lignes 2026, et une série régulière de \num{40488}\mesure{} (2012) à \num{48690}\mesure{} (2026). Aucune
source du dépôt ne donne \num{62000} ; son origine n'est pas identifiée. Le fichier brut DREES n'a pas été relu
directement (bibliothèque \texttt{openpyxl} absente).

\subsection{Limites des données}\label{don:limites}

\paragraph{Champs manquants.}
\begin{itemize}
  \item Annuaire Santé : ni âge, ni date de naissance, ni archive. 11,5~\%\mesure{} des activités dentaires
  n'ont pas de commune (0,35~\%\mesure{} pour les PH). Le sexe est la civilité déclarée. L'archive locale garde
  l'identifiant RPPS : c'est une donnée brute, non versionnée.
  \item BTS : aucun croisement secteur $\times$ âge ni secteur $\times$ sexe ; les IDE n'entrent dans le jumeau
  que sous $H_{\mathrm{ind}}$, non testée sur l'âge.
  \item CartoSanté : pas de sexe ; série SF illisible ; pas de pharmaciens ; les séries \donfic{*_flux} sont
  des pôles d'attraction des patients, pas des installations.
  \item data.ameli.fr : département seulement, sans les PH.
  \item Zonage : la date des arrêtés n'est pas servie ; deux générations seulement (« 2012 », déduite, et
  « actuelle », non datée) ; générations 2008--2011 et 2018--2023 introuvables.
\end{itemize}

\paragraph{Ruptures de série.}
\begin{itemize}
  \item MK : passage d'ADELI au RPPS en 2017 et année blanche de diplômés en 2018 \citep{d-er1159}.
  \item BTS : la feuille Lisez-moi ne signale qu'un changement du calcul de l'âge au millésime 2021. La
  « rupture de série en 2022 » écrite dans \texttt{A4} et \donfic{R/secteur} n'est pas vérifiable
  (\texttt{Audit}, C2).
  \item IFSI : changement de définition de l'abandon en 2022 et de questionnaire en 2025 \citep{d-drees491}.
  \item Substitut $I_{\mathrm{net}}$ : il ne passe le test T1 qu'à partir de 2012 ; en 2006--2011, l'écart
  atteint 63~\%\mesure{} (MK) et 201~\%\mesure{} (CD).
  \item FINESS : catégories 603 et 124 disponibles depuis 2012 et 2014 ; un transfert d'officine sous un
  nouveau numéro compte comme une fermeture et une ouverture ; Mayotte n'a pas d'officines dans FINESS avant
  2014 (écart de $-62$ à $-71$~\%\mesure{} en 2012--2013).
  \item Fichier d'inscrits 2025-26 : \num{683}\mesure{} inscrits en 2\textsuperscript{e} année d'odontologie à
  l'université Paris Cité, contre 188 à 215\mesure{} sur 2020--2024. C'est un artefact de fichier, de cause
  inconnue ; le registre l'enregistre comme une rupture de données (\donfic{R/registre}).
  \item Population : 2024--2025 reprennent 2023\suppose{} ; l'âge est interpolé entre les recensements 2012,
  2017 et 2023.
\end{itemize}

\paragraph{Secret statistique.} La DREES masque les cellules de moins de 5\mesure{} personnes ; les chargeurs ne
les comptent pas encore (reprise de l'observatoire, C6, non vérifié).

\subsection{Données traitées commitées et manifeste}\label{don:manifeste}

Les sorties de \donfic{v2/data/processed/} sont versionnées : 28\mesure{} fichiers et leur manifeste
(\texttt{git ls-files}). Le manifeste \donfic{v2/data/processed/MANIFEST.md} donne, pour chaque fichier, sa
taille en octets, les 16 premiers caractères de son SHA-256, le script qui le produit, le dernier commit de ce
script et la commande de reconstruction. Le manifeste \donfic{data/processed/MANIFEST.md} couvre les
6\mesure{} sorties de la v1 encore lues (\donfic{stock.parquet}, \donfic{population.parquet},
\donfic{stock_ameli.parquet}, \donfic{stock_commune.parquet}, \donfic{stock_nurses_salaried.parquet},
\donfic{checks.csv}).

Le contrôle se lance depuis la racine :
\begin{verbatim}
uv run python v2/data/manifest.py --check
\end{verbatim}
Il a renvoyé, le 2026-11-15, « manifestes à jour, aucun fichier des organisateurs dans la construction », avec
un code de sortie nul\mesure. Il vérifie donc deux choses : les empreintes des fichiers commités correspondent
au manifeste, et aucun fichier fourni par les organisateurs sous accord de confidentialité n'entre dans la
construction.

Deux limites :
\begin{itemize}
  \item les scripts v1 portent la mention « non commité » dans \donfic{data/processed/MANIFEST.md} : le commit
  du code qui a produit \donfic{stock.parquet} n'est pas traçable ;
  \item le manifeste prouve l'identité des fichiers, pas leur reproductibilité : reconstruire demande les
  données brutes, non versionnées, aux millésimes du tableau~\ref{don:tab-catalogue}. Les sources servies en
  continu (Annuaire Santé, CartoSanté, FINESS) ne gardent pas leurs anciens millésimes en ligne ; data.gouv
  signale d'ailleurs l'arrêt des flux FINESS utilisés, remplacés par une nouvelle version du 2026-07-20.
\end{itemize}
\endgroup
